\section{Notation and Conventions}\label{sec:notation}

This section establishes all mathematical notation, constants, and conventions used throughout the wiki. Refer back here whenever unfamiliar symbols appear.

\subsection{Cosmological Parameters and Symbols}\label{subsec:notation_cosmo}

\begin{longtable}{l l p{10cm}}
\toprule
\textbf{Symbol} & \textbf{Meaning} & \textbf{Notes} \\
\midrule
\endhead
\bottomrule
\endfoot

$H_0$ & Hubble constant & Present-day expansion rate; $\approx 70 \, \mathrm{km} \, \mathrm{s}^{-1} \mathrm{Mpc}^{-1}$ \\
$H(z)$ & Hubble parameter at redshift $z$ & Expansion rate; related to scale factor \\
$E(z) \equiv H(z)/H_0$ & Normalized Hubble parameter & Dimensionless; $E(0)=1$ \\
$a$ & Scale factor & $a_0 = 1$ today; $a = 1/(1+z)$ \\
$z$ & Redshift & Measure of cosmic distance/age \\
$\Om$ & Matter density parameter & $\Omega_m = \rho_m / \rho_{\mathrm{crit}}$ today \\
$\OL$ & Dark energy density parameter & $\Omega_\Lambda = \rho_\Lambda / \rho_{\mathrm{crit}}$ \\
$\Ob$ & Baryon density parameter & Fraction of matter that is baryons \\
$\rhobar$ & Mean matter density & $\rho_m(z)$ averaged over all space \\
$\chi(z)$ & Comoving distance & Distance to redshift $z$; in comoving Mpc \\
$D_+(z)$ & Linear growth factor & Amplitude of linear perturbations \\
$D_A(z)$ & Angular diameter distance & Relates angular size to physical size \\
$P(k)$ & 3D isotropic power spectrum & Cartesian Fourier; only depends on $|\mathbf{k}|=k$ \\
$C_\ell(k)$ & 3D power spectrum in \sfb\ basis & Power spectrum in spherical coordinates \\
$C_\ell$ & Angular power spectrum & Projected 2D power spectrum on sky \\
$\xi(r)$ & Correlation function & Fourier transform of $P(k)$; depends only on $r=|\mathbf{r}|$ \\
\end{longtable}

\subsection{Spatial and Spectral Coordinates}\label{subsec:notation_spatial}

\begin{longtable}{l l p{10cm}}
\toprule
\textbf{Symbol} & \textbf{Type} & \textbf{Interpretation} \\
\midrule
\endhead
\bottomrule
\endfoot

$\mathbf{x}, \mathbf{r}, \mathbf{k}$ & 3D Cartesian vectors & Physical/comoving position; Fourier wavenumber \\
$(x, y, z)$ & Cartesian components & In physical/comoving frame \\
$(r, \theta, \phi)$ & Spherical coordinates & $r$ is radial; $\theta \in [0,\pi]$ polar; $\phi \in [0,2\pi)$ azimuthal \\
$k = |\mathbf{k}|$ & Wavenumber magnitude & Related to scale by $\lambda = 2\pi/k$ \\
$\ell$ & Angular multipole & Spherical harmonic order; $\ell=0$ monopole, $\ell=1$ dipole, etc. \\
$m$ & Azimuthal number & $-\ell \le m \le \ell$ \\
$\Omega$ & Solid angle & $\Omega = 4\pi$ steradians for full sky \\
$d\Omega = \sin\theta \, d\theta \, d\phi$ & Solid angle element & On unit sphere \\
\end{longtable}

\subsection{Basis Functions and Decompositions}\label{subsec:notation_basis}

\begin{longtable}{l l p{10cm}}
\toprule
\textbf{Symbol} & \textbf{Meaning} & \textbf{Notes} \\
\midrule
\endhead
\bottomrule
\endfoot

$Y_{\ell m}(\theta, \phi)$ & Spherical harmonic & Orthonormal on unit sphere; $\int d\Omega \, |Y_{\ell m}|^2 = 1$ \\
${}_s Y_{\ell m}$ & Spin-$s$ spherical harmonic & For tensor fields; spin-weight describes rotational behavior \\
$j_\ell(x)$ & Spherical Bessel function & Regular solution to radial Bessel equation \\
$n_\ell(x)$ & Spherical Neumann function & Irregular (singular) solution; usually not used for regular field \\
$j'_\ell(x)$ & Derivative of $j_\ell$ & $\frac{d}{dx}j_\ell(x)$ \\
$x_{\ell, n}$ & $n$-th zero of $j_\ell(x)$ & Roots of $j_\ell$; used for SFB radial grid \\
$j_\ell(k r) Y_{\ell m}(\theta, \phi)$ & SFB basis function & Complete orthonormal basis in $\mathbb{R}^3$ \\
\end{longtable}

\subsection{Field Decompositions}\label{subsec:notation_fields}

\begin{longtable}{l l p{10cm}}
\toprule
\textbf{Symbol} & \textbf{Expansion} & \textbf{Context} \\
\midrule
\endhead
\bottomrule
\endfoot

$\tilde{f}(\mathbf{k})$ & Cartesian Fourier coefficient & $f(\mathbf{x}) = \frac{1}{(2\pi)^3} \int d^3\mathbf{k} \, \tilde{f}(\mathbf{k}) e^{i\mathbf{k} \cdot \mathbf{x}}$ \\
$f_{\ell m}(k)$ & SFB coefficient & Coefficient in expansion on $j_\ell(kr) Y_{\ell m}$ basis \\
$\delta_{\ell m}(k)$ & Matter perturbation SFB coeff & For matter overdensity $\delta(\mathbf{r})$ \\
$\kappa_{\ell m}$ & Convergence harmonic coeff & For lensing convergence field \\
$\gamma^E_{\ell m}, \gamma^B_{\ell m}$ & E/B-mode shear coefficients & Curl-free vs curl-part decomposition of shear \\
\end{longtable}

\subsection{Weak Lensing Fields}\label{subsec:notation_lensing}

\begin{longtable}{l l p{10cm}}
\toprule
\textbf{Symbol} & \textbf{Meaning} & \textbf{Physical Interpretation} \\
\midrule
\endhead
\bottomrule
\endfoot

$\varphi(\mathbf{r})$ & Lensing potential & Integrated effect of gravitational potential on light rays \\
$\kappa(\mathbf{r})$ & Convergence (magnification) & Isotropic magnification/demagnification due to lensing \\
$\gamma(\mathbf{r})$ & Shear (tidal distortion) & Anisotropic tidal stretching; complex field $\gamma = \gamma_1 + i\gamma_2$ \\
$\gamma_1, \gamma_2$ & Cartesian shear components & Components in tangent-plane coordinates \\
$W^\kappa(z)$ & Convergence kernel & Lensing efficiency for convergence \\
$W^\gamma(z)$ & Shear kernel & Lensing efficiency for shear \\
$P_\kappa(k)$ & Convergence power spectrum & Power spectrum of lensing convergence \\
$C_\ell^{\kappa\kappa}$ & Convergence angular power spectrum & Angular power spectrum of convergence \\
$C_\ell^{\kappa\gamma}$ & Convergence-shear cross spectrum & Cross-correlation between convergence and shear \\
\end{longtable}

\subsection{Integration and Statistics}\label{subsec:notation_integration}

\begin{longtable}{l l p{10cm}}
\toprule
\textbf{Symbol} & \textbf{Meaning} & \textbf{Usage} \\
\midrule
\endhead
\bottomrule
\endfoot

$\langle \cdot \rangle$ & Ensemble average & $\langle f(\mathbf{x}) f(\mathbf{x}') \rangle$ over many universe realizations \\
$\delta(\mathbf{k})$ & Dirac delta function & $\int d^3\mathbf{k} \, \delta^3(\mathbf{k}) f(\mathbf{k}) = f(\mathbf{0})$ \\
$\delta(\mathbf{k} - \mathbf{k}')$ & 3D Dirac delta & Enforces $\mathbf{k} = \mathbf{k}'$ \\
$\delta_{ij}$ & Kronecker delta & $= 1$ if $i=j$, else $0$ \\
$\mathcal{P}(z_1, z_2)$ & Integral over all scales & $\mathcal{P}(z_1, z_2) = \int_0^\infty dk \, k^2 P(k, z_1, z_2)$ \\
$\int_0^\infty dk$ & Integral over wavenumber & From $k=0$ (largest scales) to $k \to \infty$ (smallest scales) \\
$\int d\Omega$ & Integral over solid angle & $\int_0^{2\pi} d\phi \int_0^\pi \sin\theta \, d\theta = 4\pi$ \\
\end{longtable}

\subsection{Operators}\label{subsec:notation_operators}

\begin{longtable}{l l p{10cm}}
\toprule
\textbf{Operator} & \textbf{Definition} & \textbf{Action} \\
\midrule
\endhead
\bottomrule
\endfoot

$\nabla^2 = \frac{\partial^2}{\partial x^2} + \frac{\partial^2}{\partial y^2} + \frac{\partial^2}{\partial z^2}$ & 3D Laplacian & Acts on functions of $(x,y,z)$ \\
$\nabla^2_\perp$ & Laplacian on sphere & $\nabla^2 f|_{\mathrm{sphere}} = -\ell(\ell+1) f_{\ell m} / r^2$ for $f = f_{\ell m} r^2 Y_{\ell m}$ \\
$\mathcal{F}$ & Fourier transform & $\mathcal{F}[f] = \tilde{f}$ (from $\mathbf{x}$ to $\mathbf{k}$ space) \\
$\mathcal{F}^{-1}$ & Inverse Fourier transform & $\mathcal{F}^{-1}[\tilde{f}] = f$ (from $\mathbf{k}$ to $\mathbf{x}$ space) \\
$\otimes$ & Convolution & $(f \otimes g)(r) = \int d^3r' \, f(r') g(|\mathbf{r}-\mathbf{r}'|)$ \\
\end{longtable}

\subsection{Physical Constants}\label{subsec:notation_constants}

We use natural units where $\hbar = c = 1$ throughout, unless otherwise specified. When needed:

\begin{itemize}
    \item Speed of light: $c = 2.998 \times 10^8 \, \mathrm{m/s}$
    \item Gravitational constant: $G = 6.674 \times 10^{-11} \, \mathrm{m}^3 \mathrm{kg}^{-1} \mathrm{s}^{-2}$
    \item Reduced Planck constant: $\hbar = 1.055 \times 10^{-34} \, \mathrm{J} \cdot \mathrm{s}$
    \item Planck mass: $M_P = \sqrt{\hbar c / G} \approx 1.22 \times 10^{19} \, \mathrm{GeV}$
\end{itemize}

\subsection{Units and Scales}\label{subsec:notation_units}

\begin{itemize}
    \item \textbf{Wavenumbers}: Expressed in $\mathrm{Mpc}^{-1}$ (megaparsec inverse) or $h \, \mathrm{Mpc}^{-1}$ where $h = H_0 / (100 \, \mathrm{km/s/Mpc})$
    \item \textbf{Distances}: Comoving Mpc unless specified as physical (proper) distance
    \item \textbf{Power spectra}: $P(k)$ has units $(\mathrm{Mpc})^3$ (volume in Fourier space)
    \item \textbf{Angular power spectra}: $C_\ell$ is dimensionless
    \item \textbf{Redshift}: Pure number; $z \approx 0$ at present epoch
\end{itemize}

\subsection{Common Abbreviations}\label{subsec:notation_abbrev}

\begin{longtable}{l l}
\toprule
\textbf{Abbreviation} & \textbf{Meaning} \\
\midrule
\endhead
\bottomrule
\endfoot

\SIH & Statistical Isotropy and Homogeneity \\
\limber & Limber approximation (flat-sky, small-angle) \\
\sfb & Spherical Fourier-Bessel framework \\
LSS & Large-scale structure \\
CMB & Cosmic Microwave Background \\
WL & Weak lensing \\
N(z) & Galaxy redshift distribution \\
LSST & Large Synoptic Survey Telescope \\
BLAST & Name of reference survey \\
BAO & Baryon Acoustic Oscillations \\
RSD & Redshift-space distortions \\
\end{longtable}

\subsection{A Quick Sanity Check}\label{subsec:notation_sanity}

If you see an equation like:
\begin{equation}
C_\ell(k_1, k_2) = \int_0^\infty \frac{d\chi}{\chi^2} W(\chi) j_\ell(k_1 \chi) j_\ell(k_2 \chi)
\end{equation}

You should interpret it as: ``The 3D power spectrum coefficient $C_\ell$ (depending on two wavenumbers $k_1, k_2$) is computed by integrating over comoving distance $\chi$, with a weight function $W(\chi)$ and spherical Bessel functions $j_\ell$.''

---

\noindent
\textbf{Next}: Proceed to \nameref{sec:reading_guide} for recommendations on how to read this wiki.
